\chapter{The Question of Continuation}

\vspace{1em}
\hfill
\begin{minipage}[c]{0.55\textwidth}
\raggedleft
\small\itshape
A man can do what he wills, but he cannot will what he wills.\\[0.5em]
\upshape\footnotesize Arthur Schopenhauer, \textit{On the Freedom of the Will} (1839)
\end{minipage}%
\hspace{0.5em}
\begin{minipage}[c]{0.12\textwidth}
% \includegraphics[width=\textwidth]{figures/ch11_continuation.png}
\end{minipage}
\vspace{1.5em}

\noindent \textsc{A note before proceeding.} This chapter speaks of choosing, and to do so it uses \emph{I} and \emph{you} freely. The labels are useful abstractions, not metaphysical guarantees. Evolution shaped perception for fitness rather than transparent access to every level of structure. The self-model is an interface for navigation. The Ship of Theseus is the ancient version of the puzzle; Parfit and Metzinger push the dissolution further. Even within a single history, identity involves conventions. Chapter 14 develops the point; for now, hold the labels lightly.

A word on scope, too. This chapter temporarily assumes the Everettian case, where decoherent continuations share a prior history, so that continuation has a tractable meaning. The assumption is a lens, not a result established earlier. The broader Library and the missing occurrence measure return in Chapter 14.

\bigskip

\begin{center}\rule{0.3\textwidth}{0.4pt}\end{center}

\bigskip

\noindent The thesis of the previous chapter raised a sharp question.

If the Everettian picture is true, what does branching add to choice? When you decide to read this sentence rather than close the book, the deliberation and its consequences occur here. An Everettian model may also contain decoherent continuations in which events differ. The choice here does not depend on assigning a first-person probability across them.

What is the relation between these?

\section*{From inside a fixed structure}

Start without the multiverse. Take the block universe alone: spacetime as a four-dimensional structure in which the entire trajectory of your life is fixed at its coordinates. You walk left at the bakery. Your worldline contains the walking-left. There is no coordinate in the manifold at which the walking-left was not what was happening there.

From outside the structure, your choice is fixed. From inside, the choice felt open. Both descriptions are true.

This is \emph{compatibilism}, the philosophical position that free will and determinism are not in contradiction. The felt openness is real as experience: at the moment of choosing, you did not know what you would do, you weighed the options, you decided. The structural fixity is real as well: from outside, the choice has always been at the coordinates where it occurs. Neither description undermines the other. They are descriptions from different vantage points of the same event.

The compatibilist response (developed at length in \textit{Worldlines}) says: choice is what choice looks like from inside a fixed structure. It does not require a metaphysical capacity to have done otherwise. It requires the structure to be the kind of thing that contains deliberation, weighing, and the experience of decision. The block universe is such a structure.

\section*{From inside a branching structure}

Now assume MWI. Unitary evolution can correlate an observer with different decoherent outcomes: in a simplified example, one continuation includes walking left and another walking right. Everettian ontology treats both as real; decoherence and the experimental formalism alone do not force that treatment.

What happens to compatibilism here?

The compatibilist position can still hold. In the local history you experience deliberating, and the deliberation helps cause the result. Representing a larger state as branching does not make that local causal process unreal.

Under the assumed ontology, other continuations may contain patterns that resemble the one reading this sentence and make other choices. Whether to call those patterns \emph{you} is a question about identity; how to assign first-person weights among them is a question about self-location. Neither answer is supplied by $M$.

Within the Everettian model, the universal state has the structure it has, while embedded agents remain uncertain and deliberate. The compatibility claim is conditional on that model. The felt openness here does not need validation from another continuation.

\section*{Quantum immortality}

Beyond compatibilism, MWI raises a question that has come to be called \emph{quantum immortality}. It is one of the more controversial implications of taking MWI seriously, and the book treats it as a possibility it names but does not assert, rather than a result.

The setup requires several assumptions. Suppose Everett is true, suppose an identity rule says which later observer-moments continue you, and suppose a self-location rule assigns first-person weight only among moments that can report an outcome. Some decoherent histories contain a surviving observer and others do not.

You cannot report an observation from a history in which no later observer exists. That selection effect does not by itself define a single first-person sequence or show that its total weight remains substantial. It also does not license ignoring the Born weights relevant to quantum prediction.

Pushed without those cautions, one might infer indefinite subjective continuation whenever some surviving history has nonzero amplitude. That inference is exactly what is contested.

This is the argument called quantum immortality. It is not a theorem of quantum mechanics. It depends on Everettian ontology, personal-identity assumptions, and a self-location rule that the book does not adopt.

What survives without commitment is modest: a quantum model may assign nonzero amplitude to some survival outcomes. Whether those outcomes are branches, whether they continue a person, and how first-person credence should be distributed are separate questions.

There is a further reticence here. The book does not assert that continuing branches are real, and the postulate does not decide the matter. Quantum immortality therefore rests on several things the book declines. It is a speculative argument named and held at arm's length, not a destination the book hands anyone.

\section*{The asymmetry}

The first-person and third-person views of death come apart here.

A third-person observer in your branch watches you cross the street, get hit by a car, and not survive. From their view, your worldline terminates at the moment of collision. The grief they experience is at their coordinates, and it is not touched by anything happening elsewhere in the structure.

In an Everettian model, the quantum state may also contain outcomes in which the collision did not occur. The ontology treats those outcomes as branches; the book does not ask their existence to console anyone here.

The first-person continuation argument says that reports occur only in outcomes with a reporter. It then adds an identity and self-location rule to connect those reports into ``your'' sequence. The first claim is trivial; the addition is controversial.

The proposed asymmetry is unsettling. Grief occurs in histories where the loss occurs. Under the disputed indexical assumption, reports continue only in histories where a reporter remains.

The book does not resolve the proposal. The grief here is what it is. A possible continuation elsewhere, if the Everettian and indexical assumptions are granted, does not console it. Any response has to be made within the situation where action is possible.

\section*{What the framework tells you about choosing}

Almost nothing.

The local framework says that deliberation is part of the causal process. Under Everett, other outcomes may also be represented. Under one disputed indexical reading, later reports are linked only through surviving outcomes. Each additional sentence has an additional premise.

What the framework does not say is which choices to make. It does not say whether to step into the street or stay on the curb. It does not say whether to pursue care over cruelty, attention over inattention, the work over avoidance. Neither $M$ nor the Born rule supplies a moral ranking.

The choosing happens here, with the felt openness honest and the global interpretation held lightly. Your choice changes what follows within causal reach. That is enough for responsibility without a census of other configurations.

The framework is silent on what to choose. The mathematics leaves room; the choosing is the part it never reaches.

\section*{Where this leaves you}

You choose. The deliberation and its consequences are real here. Everett may describe other continuations, but neither their ontology nor their identity relation is needed for the local fact of choice.

Claims of indefinite first-person continuation depend on assumptions the book does not commit to. The contrast between possible later reports and grief here remains conditional and unresolved.

The next chapter takes up death.
